\documentclass[9pt,a4paper]{extarticle}

\usepackage[utf8]{inputenc}
\usepackage{geometry}
\usepackage{xcolor}
\usepackage{hyperref}

\geometry{
	a4paper,
	left=0.25cm,
	right=0.25cm,
	top=0.25cm,
	bottom=0.25cm
}

\hypersetup{
	colorlinks=true,
	linkcolor=blue,
	urlcolor=blue,
	pdfborder={0 0 0}
}

\renewcommand{\section}[1]{%
	\vspace{2pt}%
	{\large\bfseries #1}\par\nobreak\vspace{-0.7em}%
	\rule{\linewidth}{0.4pt}\par\nobreak%
	\vspace{-2pt}%
}

\newenvironment{tightlist}{%
	\vspace{-3pt}%
	\begin{itemize}%
		\setlength{\leftmargin}{0pt}%
		\setlength{\itemindent}{0pt}%
		\setlength{\itemsep}{-2pt}%
		\setlength{\parsep}{0pt}%
		\setlength{\topsep}{0pt}%
		\setlength{\partopsep}{0pt}%
	}{%
	\end{itemize}%
	\vspace{-6pt}%
}
\setlength{\parskip}{0pt}
\setlength{\parindent}{0pt}


\begin{document}
	\begin{center}
		{\huge\textbf{Cameron Badman}} \\[0.3em]
		\href{mailto:cbadwork@gmail.com}{cbadwork@gmail.com} |
		+61 476056341 |
		\href{https://github.com/CameronBadman}{GitHub}|
		\href{https://www.linkedin.com/in/cameron-badman-5314ba1b8/}{LinkedIn}
	\end{center}
	\section{Professional Experience}
	\textbf{Software Engineer, Gumnut.dev} \hfill Sep 2025--Feb 2026
	\begin{tightlist}
		\item Developed advanced Zustand and Node.js integrations for real-time collaborative state management
		\item Implemented Go data structures and sub-millisecond optimisations for text processing and editing pipelines
		\item Researched and applied garbage collection strategies in distributed computing environments
		\item Built presence cursor system enabling real-time multi-user awareness across collaborative sessions
		\item \textit{Technologies: Go, Node.js, TypeScript, Zustand, WebSockets, Data Structures}
	\end{tightlist}
	\textbf{Software Engineer, Base} \hfill Dec 2024--Sep 2025
	\begin{tightlist}
		\item Implemented Github CI/CD pipelines to streamline development and deployment processes
		\item Managed a small group of international developers across multiple time zones
		\item Performed significant refactors to modularise code, enabling easy switching between monolithic and services-based architecture
		\item Migrated from Node PayPal SDK v0.8 to 1.1.0, improving payment processing reliability
		\item \textit{Technologies: PostgreSQL, Node.js, React Native}
	\end{tightlist}
	\textbf{DevOps Engineer, Cyberloop (Oil tech)} \hfill Jan 2023--Feb 2024
	\begin{tightlist}
		\item Implemented CI/CD pipelines in bitbucket that reduced manual testing time by 75\%
		\item Refactored Cypress tests and wrote majority of tests for multiple complex SaaS products
		\item Reduced migration costs by 20\% by creating bash script to compress unstructured data on Firebase
		\item Developed GenAI-powered tool that assigns labels for difficulty and time automatically
		\item Maintained and deployed Docker containers using Portainer and Helm on Google Cloud Engine
		\item \textit{Technologies: GCP, Docker, BitBucket pipelines, Jest, Playwright}
	\end{tightlist}
	\textbf{CSSE6400 Tutor, University of Queensland} \hfill Feb 2026--Present
	\begin{tightlist}
		\item Teaching cloud engineering concepts including Terraform, AWS, and Docker to undergraduate students
		\item Write assessment solutions and testing suites for course assignments
	\end{tightlist}
	\textbf{Python Tutor, Junior Engineers} \hfill Jun 2021--Present
	\begin{tightlist}
		\item Rewrote programming classes for children and teens, making complex coding concepts easier to understand
		\item Led advanced classes for 8-16 year olds, teaching Python OOP, Scratch, Arduino, Minitendo
	\end{tightlist}
	\section{Education}
	\textbf{University of Queensland -- Bachelors of Computer Science/Commerce} \hfill 2021--Present
	\par\vspace{1pt}
	\section{Projects - \textbf{All projects available on} \href{https://github.com/CameronBadman}{GitHub}}
	\textbf{Hippocampus}: \textit{File-based vector database built for AI agent memory} \href{https://github.com/CameronBadman/Hippocampus}{[GitHub]}
	\begin{tightlist}
		\item Built custom O(512 log n) binary search achieving 4,200x faster exact search than FAISS at 10k vectors
		\item Designed file-based architecture with per-agent isolated .bin files and binary serialization (2KB per node)
		\item Deployed serverless infrastructure on AWS Lambda, EFS, and S3 using Terraform
		\item Implemented agent-to-agent orchestration with AWS Bedrock Nova for automated memory curation
		\item \textit{Technologies: Go, AWS Lambda, Bedrock, EFS, S3, Terraform, Python}
	\end{tightlist}
	\textbf{IT Support AI Agent}: \textit{Self-evolving IT support agent with tiered semantic memory} \href{https://github.com/CameronBadman/IT-Agent-Thing}{[GitHub]}
	\begin{tightlist}
		\item Architected three siloed Hippocampus instances for company, group, and user-level policy memory
		\item Agent queries memories to contextualise responses, e.g. never returning Windows commands for a Mac-only organisation
		\item Knowledge base evolves autonomously over time as policies and preferences are learned and stored
		\item Built for the Claude Code Hackathon using Ollama Mistral 7B with structured tool calling
		\item \textit{Technologies: Python, Go, Ollama, Hippocampus, Docker}
	\end{tightlist}
	\textbf{Gleam DynamoDB Client}: \textit{First Gleam DynamoDB client} \href{https://github.com/CameronBadman/DynamoDb-gleam-client}{[GitHub]}
	\begin{tightlist}
		\item Implemented and interfaced with the low-level AWS API integrations and documentation
		\item Designed type-safe monadic recursive dynamo json framework
		\item Learnt the gleam language in spite of documentation and llm support being non-existent
		\item Integrated AWS signature authentication and secure communication protocols
		\item \textit{Technologies: Gleam, AWS DynamoDB, AWS SDK, Functional Programming}
	\end{tightlist}
	\textbf{Collaborative Drawing WebApp}: \textit{Real-time collaborative drawing application} \href{https://github.com/CameronBadman/Canvis-collab-webapp}{[GitHub]}
	\begin{tightlist}
		\item Developed custom HTML canvas drawing system with real-time synchronization
		\item Built Go microservices with goroutines and WebSocket connections for real-time collaboration
		\item Deployed on AWS using containerized architecture with Kubernetes orchestration
		\item Implemented conflict resolution algorithms for concurrent drawing operations
		\item \textit{Technologies: Go, JavaScript, Docker, Kubernetes, AWS, WebSockets}
	\end{tightlist}
	\section{Open-Source Contributions}
	\textbf{goproxy}: \textit{Go HTTP proxy library forked by Stripe, Minikube, and Grafana} \href{https://github.com/elazarl/goproxy}{[GitHub]}
	\begin{tightlist}
		\item Developed HTTP Archive (HAR) middleware with Prometheus/Grafana integration using goroutines
		\item Modernizing project's testing infrastructure, transitioning to modern testing practices using testify
	\end{tightlist}
	\section{Community}
	\textbf{Startup Social Brisbane} -- Host \hfill \href{https://www.meetup.com}{Meetup}\\
	\textbf{IT Social Brisbane} -- Host \hfill \href{https://www.meetup.com}{Meetup}
	\par\vspace{1pt}
	\section{Technologies and Skills}
	\textbf{Languages:} Go, Java, Python, TypeScript, Bash, SQL, Gleam, C, C++
	
	\textbf{Web Technologies:} React, Vue.js, ASP.NET, HTML, CSS
	
	\textbf{Cloud \& DevOps:} Amazon Web Services (AWS), Google Cloud Platform (GCP), Docker, Kubernetes, Terraform, Nix
	
	\textbf{Databases:} PostgreSQL, SQLite, Redis, Cassandra, DynamoDB, Valkey
	
	\textbf{Tools:} Git, GitHub, Playwright, CI/CD, Kafka, Atlassian
	
	\textbf{Operating Systems:} Windows, Mac, Linux, Arch-Linux, NixOS, Ubuntu, Gentoo
	
\end{document}